\subsection*{真实模型加载闭环：Tokenizer、权重格式与首 token trace}

\section{为什么真实模型加载是第三册的分水岭}

第三册已经有 tensor、KV cache、量化、IR、runtime 和 backend 账本。但如果模型资产仍然停留在手写 manifest 或 tiny fixture，读者还没有真正进入 AI Infra。真实推理系统的第一道工业边界，是把磁盘上的 tokenizer、config、权重分片、量化 metadata 和 chat template 变成一个可验证的 executable plan。任何一个字段错了，后面的 kernel 再快也只是更快地产生错误 logits。

本章目标是建立“真实模型加载闭环”的验收标准：

\begin{lstlisting}[numbers=none,caption={真实模型加载闭环}]
model directory or package
  -> tokenizer contract
  -> config and architecture contract
  -> tensor table and weight shards
  -> quant descriptor
  -> name mapping and shape verifier
  -> typed graph
  -> memory map or staged load
  -> first token trace
  -> logits oracle
\end{lstlisting}

完成这条链路以后，推理系统才从原理说明进入可训练、可复查的工程闭环。

\section{Tokenizer 是模型合同的一部分}

Tokenizer 不是服务边缘的小工具。它决定 token id、special token、BOS/EOS、padding、chat template、位置编号和 embedding 行索引。tokenizer 版本变化，可能让同一个 prompt 变成不同 token 序列；chat template 变化，可能让角色、换行和特殊标记不同；special token id 错，可能直接访问错误 embedding 行。

\begin{center}
\small
\begin{tabular}{p{0.24\linewidth}p{0.34\linewidth}p{0.30\linewidth}}
\toprule
字段 & 必须验证什么 & 错误后果 \\
\midrule
vocab size & embedding/lm head 行数匹配 & 越界或 logits 维度错 \\
token id mapping & special token id 稳定 & BOS/EOS/UNK 错误 \\
normalizer & 文本归一化规则 & 同一文本 token 不同 \\
pre-tokenizer & 空格、字节、unicode 处理 & prompt 重放失败 \\
chat template & role 到文本的展开规则 & 对话语义漂移 \\
padding side & position id 和 mask & RoPE/KV 位置错 \\
\bottomrule
\end{tabular}
\end{center}

报告不能只写“加载 tokenizer 成功”。它要固定一个 prompt fixture，输出 token ids、special token 位置、chat template 展开文本 hash 和 tokenizer 配置 hash。若 tokenizer.json、SentencePiece model 或 GGUF 内嵌 tokenizer 互相冲突，verifier 必须拒绝或明确选择优先级。

\subsection{当前 lab 的 tokenizer 合同}

当前 \filepath{books/cpu-volume-3/labs/linux\_cpu\_inference} 已经把 tokenizer 从“正文要求”推进到可运行 fixture。它不是完整 BPE 或 SentencePiece 实现，而是一个最小合同：BOS/EOS、UNK 策略、chat template hash、词表、prompt fixture 和合同 hash。

\begin{lstlisting}[numbers=none,caption={tiny tokenizer fixture 的核心字段}]
LCQI_TOKENIZER_V1
bos_id 0
eos_id 4
unk_id -1
chat_template_hash tiny-chat-template-v1
vocab_size 5
vocab:
  0 <bos>
  1 tok1
  2 tok2
  3 tok3
  4 <eos>
\end{lstlisting}

运行 \code{lcqi\_trace} 时，trace 会先输出完整 tokenizer 结果，再把 decoder 输入单独传给 reference decoder：

\begin{lstlisting}[numbers=none,caption={tokenizer trace 的最小硬证据}]
tokenizer,0,-1,5,1,i32,contiguous,10,4,0;1;2;3;4
tokenizer_contract,0,-1,scalar,scalar,u64,fnv1a,13452902845388333734,0,tiny-chat-template-v1
\end{lstlisting}

这个设计故意把两层分开。Tokenizer 合同包含 BOS/EOS、template 和词表身份；decoder 合同只接收本次 tiny 模型需要的正文 token。真实 LLaMA/Qwen 类模型也要做同样分层：chat template、special token、padding、position id 和 decoder input ids 都要在 trace 中可见。否则换了 tokenizer 或 template 后，后端 kernel 仍然能运行，却已经不再处理同一个请求。

下一步从 fixture 走向真实格式时，不应删除这个合同，而是扩展它：

\begin{lstlisting}[numbers=none,caption={从 fixture 扩展到真实 tokenizer 的门禁}]
tokenizer_gate:
  tokenizer_source: tokenizer.json / sentencepiece.model / gguf metadata
  version_or_hash:
  normalizer:
  pre_tokenizer:
  special_tokens:
  chat_template_hash:
  prompt_fixture_ids:
  detokenize_roundtrip:
  embedding_vocab_match:
\end{lstlisting}

完整 BPE/SentencePiece 的复杂度在分词算法本身；工程风险在合同漂移。实验代码先固定合同，是为了让后续替换真实 tokenizer 时有可回归的出口。

\section{GGUF、safetensors 和分片权重}

GGUF 更像本地推理包：metadata、tokenizer、tensor table 和量化类型常放在同一文件中。safetensors 更像安全 tensor 容器：header 描述 dtype、shape 和 data offsets，权重可能分片，config 和 tokenizer 通常在旁边文件中。两者都不是“读一堆 float”这么简单。

\begin{lstlisting}[numbers=none,caption={模型资产清单}]
model_asset:
  format: gguf / safetensors / custom_fixture
  architecture: llama_like / qwen_like / other
  config:
    layers
    hidden
    intermediate
    attention_heads
    kv_heads
    rope_theta
    vocab_size
  tokenizer:
    source
    vocab_hash
    chat_template_hash
  tensors:
    name
    dtype
    shape
    data_offset
    shard_id
    quant_type
    checksum
\end{lstlisting}

真实加载器最容易错的是名字映射。不同生态对同一权重可能命名不同：\code{layers.0.attention.wq.weight}、\code{model.layers.0.self\_attn.q\_proj.weight}、\code{blk.0.attn\_q.weight}。加载器不能靠字符串猜完就算成功；它要把每个外部 tensor 映射到内部 graph value，并记录 unmapped、duplicate、shape mismatch、dtype mismatch 和 transpose requirement。

\subsection{当前 lab 的 safetensors manifest gate}

当前 lab 已经补了一个很小的 safetensors header 解析器。它不声称能加载完整 LLaMA/Qwen 权重，也不处理所有 JSON 细节；它只把真实模型格式最关键的目录事实落到可测试代码：header 长度、\code{\_\_metadata\_\_}、tensor name、dtype、shape、\code{data\_offsets}、payload 边界，以及 dtype/shape 推导出的字节数是否等于 offset 区间。

\begin{lstlisting}[numbers=none,caption={safetensors manifest gate 的最小证据}]
SafeTensorManifest:
  header_size
  data_start_offset
  metadata:
    format -> lcqi-fixture
  tensors:
    model.embed_tokens.weight:
      dtype = F32
      shape = [2, 3]
      data_offsets = [0, 24]
    model.layers.0.attn.q_proj.weight:
      dtype = F16
      shape = [4, 3]
      data_offsets = [24, 48]
\end{lstlisting}

测试会运行时生成一个合法 fixture，并验证两个 tensor 的名称、dtype、shape 和 offset；还会生成 \code{data\_offsets} 超出 payload、dtype/shape 字节数不匹配的坏文件，要求加载器拒绝。这个 gate 的意义是把“权重格式导入”从口号压到第一层工程边界：任何后续 name mapping、shape verifier、mmap 或分片加载，都必须先建立可信 tensor table。若连 offset 范围和字节数都不检查，后端 kernel 读到的只是任意字节流。

下一阶段应在这个 gate 上继续扩展，而不是绕开它：

\begin{lstlisting}[numbers=none,caption={从 safetensors header 到真实加载器的下一层}]
safetensors_manifest
  -> shard index and tensor table merge
  -> external name to canonical role mapping
  -> config-derived expected shape
  -> dtype and quant descriptor verification
  -> mmap/read policy
  -> first token trace
\end{lstlisting}

这样写能保持边界诚实：当前 lab 已经验证真实格式的目录合同，但还没有完成完整权重解释和首 token 真实模型 oracle。

\subsection{真实 small 模型 smoke：先证明开源资产能在本机跑}

只靠 tiny fixture 训练正确性还不够；读者还需要见到一个真实开源模型文件怎样被本地推理栈加载。第三册 lab 因此新增了一个显式 smoke，而不是把它藏进普通 CTest：

\begin{lstlisting}[numbers=none,caption={运行 SmolLM2 small 开源模型 smoke}]
make cpu3-smollm2-smoke
\end{lstlisting}

这个 smoke 固定模型为 \code{HuggingFaceTB/SmolLM2-135M-Instruct}，使用 \code{bartowski/SmolLM2-135M-Instruct-GGUF} 中的 \code{SmolLM2-135M-Instruct-Q4\_K\_M.gguf}。它是 135M 参数级别的 small 模型，不是 tiny 随机夹具。脚本会校验文件大小 \code{105454432} 字节和 SHA256：

\begin{lstlisting}[numbers=none,caption={SmolLM2 small GGUF 身份}]
2e8040ceae7815abe0dcb3540b9995eaa1fa0d2ca9e797d0a635ae4433c68c2d
\end{lstlisting}

脚本还会拉取固定 commit 的 \code{llama.cpp}，构建 \code{llama-simple-chat}，用 \code{-ngl 0} 强制 CPU 路径，输入短 prompt，并把模型 identity、命令、退出码、assistant response 和运行输出尾部写入 \filepath{books/cpu-volume-3/results/lcqi-smollm2-small-model-smoke.txt}。

这条 smoke 的结论必须写窄。它证明真实 GGUF small 模型能在当前机器上完成 tokenizer/chat template、prefill、decode 和文本输出；它不证明 LCQI 自研引擎已经支持 GGUF，也不证明模型回答质量正确。tiny fixture 继续负责逐层 golden trace；small 模型 smoke 负责把“真实开源模型确实跑过”变成可复查证据。

\section{shape verifier 要从 config 推导，而不是照抄文件}

验证器不能只检查“文件里写了 shape”。它要从模型 config 推导每个 tensor 的期望 shape，再和文件形状比对。例如 decoder-only GQA 模型中：

\begin{lstlisting}[numbers=none,caption={LLaMA-like shape 推导}]
hidden = H
intermediate = I
layers = L
query_heads = QH
kv_heads = KVH
head_dim = H / QH

tok_embeddings.weight: [vocab, H]
layers.i.attn.wq.weight: [QH * head_dim, H]
layers.i.attn.wk.weight: [KVH * head_dim, H]
layers.i.attn.wv.weight: [KVH * head_dim, H]
layers.i.attn.wo.weight: [H, H]
layers.i.ffn.w1.weight: [I, H]
layers.i.ffn.w2.weight: [H, I]
layers.i.ffn.w3.weight: [I, H]
output.weight: [vocab, H]
\end{lstlisting}

如果 \code{H \% QH != 0}，\code{head\_dim} 不合法；如果 \code{KVH} 不能整除或被 \code{QH} 映射，GQA 关系不合法；如果 vocab 与 tokenizer 不一致，embedding 和 logits 语义不合法。verifier 应输出诊断，而不是让后端 kernel 在运行时越界。

\section{mmap、lazy load 和 prepare 的边界}

真实模型通常很大。加载策略至少有三类：一次性读入内存，mmap 文件并按需触页，mmap 后 prepare 阶段主动触碰和打包。mmap 成功不代表权重已经进入内存；第一次 decode 若触发大量缺页，会把 cold fault 混进 TPOT。prepare 阶段若做 packed weight、NUMA first-touch、checksum 和量化 metadata 解析，也必须和 request latency 分开报告。

\begin{lstlisting}[numbers=none,caption={加载阶段边界}]
OpenPackage
  -> ParseMetadata
  -> VerifyTokenizerAndConfig
  -> BuildTensorTable
  -> MapOrReadWeights
  -> VerifyShapesAndDTypes
  -> BuildTypedGraph
  -> PreparePackedWeights
  -> WarmPagesIfRequired
  -> ReadyForPrefill
\end{lstlisting}

服务报告要区分 load time、prepare time、first prefill time、steady decode time。若把 prepare 放进第一个请求，TTFT 会很高；若把 prepare 从报告中删掉，系统容量和冷启动成本又被隐藏。正确做法是两个数字都保留。

\section{首 token trace：加载闭环的最终证据}

真实加载闭环的第一份硬证据，不是回答自然语言问题，而是固定 prompt 的首 token trace。它应该记录每层关键中间值的 shape、dtype、layout、checksum 或小片段，以及最终 logits/top-k。trace 不需要保存巨大 tensor，但必须足以定位第一个偏离点。

\begin{lstlisting}[numbers=none,caption={first token trace 字段}]
first_token_trace:
  model_id
  asset_hashes
  tokenizer_hash
  config_hash
  prompt_text_hash
  token_ids
  backend_plan
  for each layer:
    rmsnorm_checksum
    qkv_shape_dtype_layout
    rope_position
    kv_append_slot
    attention_score_summary
    mlp_checksum
  logits:
    shape
    dtype
    checksum
    top_k
  sampler:
    temperature
    seed
    selected_token
\end{lstlisting}

如果首 token trace 能稳定复现，后续优化才有基准。换 tokenizer、换权重格式、换量化 profile、换 backend，都应先对齐这个 trace。否则“生成结果看起来差不多”无法支撑工程判断。

\section{错误案例库}

\begin{center}
\small
\begin{tabular}{p{0.26\linewidth}p{0.34\linewidth}p{0.28\linewidth}}
\toprule
错误 & 现象 & 应有诊断 \\
\midrule
tokenizer vocab 与 embedding 不一致 & 访问越界或 logits 维度错 & 报告两者来源和实际值 \\
chat template 缺失 & 对话 prompt 语义不同 & 报告模板 hash 和展开文本 hash \\
Q/K/V shape 映射错 & shape 可能仍能乘，但 head 语义错 & 指出 layer、tensor、期望 shape \\
RoPE theta 缺失或默认错 & 短 prompt 近似，长上下文漂移 & 报告 config 来源和默认策略 \\
weight transpose 方向错 & logits 大幅偏离 & name mapping 中记录 layout 转换 \\
量化 group size 错 & 输出可运行但误差大 & QuantDesc scale count 失败 \\
mmap 冷触页混入请求 & 首次 TTFT 异常高 & load/prepare/request 分段报告 \\
\bottomrule
\end{tabular}
\end{center}

\section{验收标准}

第三册的真实模型加载闭环至少要交付：

\begin{enumerate}
  \item tokenizer fixture：文本、展开文本 hash、token ids、special token 位置。
  \item asset manifest：config、tensor table、dtype、shape、offset、checksum。
  \item name mapping report：外部 tensor 到内部 graph value 的映射。
  \item shape verifier：正例通过，坏例准确拒绝。
  \item load/prepare/request 分段时间。
  \item first token trace：中间 checkpoint、logits top-k 和 sampler 状态。
  \item 回归门禁：tokenizer、config、权重、量化和 backend 任一变化都触发 trace 对比。
\end{enumerate}

\begin{keyidea}
真实模型加载不是附属功能，而是 AI Infra 的入口合同。没有 tokenizer、权重格式、shape verifier 和首 token trace，后端优化没有可信基线。
\end{keyidea}
